\section{Recursive aggregation}\label{sec:recursion}

The implemented N-to-1 aggregation runs the verifier of \S\ref{sec:e2e-unrolled} \emph{inside} the VM: one compiled guest replays the transcripts, sumchecks, ring switches, and stacked Ligerito openings of $n_{\mathrm{rec}}\ge1$ ordinary VM proofs for a fixed inner program. The proofs may have different announced table and commitment sizes; the guest dispatches among the certified shapes without recompilation. The VM is well suited to this because $\E$ is its machine-word field and its Fiat--Shamir hash is the \texttt{BLAKE3} opcode. One class of operations, however, does not fit.

\subsection{Deferred evaluation claims}\label{sec:deferred-claims}

Twice during verification, a large \emph{fixed} multilinear must be evaluated at a random point:
\begin{itemize}
\item \textbf{the bytecode}: during leaf decomposition (\S\ref{sec:e2e-unrolled}, \emph{Bus}), the verifier evaluates the nine public program columns, padded to sixteen selector slots under four selector variables, in one pass over the whole program;
\item \textbf{the flock matrices}: the \texttt{BLAKE3} relation is proven by flock (\S\ref{sec:tab-blake3}), whose lincheck ends by evaluating $\widetilde A_0$ and $\widetilde B_0$, the per-block R1CS matrices: $28$-variable multilinears with $\approx 21\mathrm{M}$ nonzero Boolean entries between them.
\end{itemize}
Either evaluation is one pass over a fixed object. This takes milliseconds natively, but in-circuit it would dwarf the entire remaining verifier, in cycles and in bytecode size alike.

The aggregation guest therefore does not evaluate these fixed tables. It verifies every surrounding reduction and exports only the resulting points and claimed values.

There are three fixed polynomials: the bytecode and the two matrices $A_0,B_0$. Verifying each ordinary sub-proof produces one claim on each, for $3n_{\mathrm{rec}}$ claims. A fresh Fiat--Shamir transcript absorbs all sub-statements and claim data, then batches them with two sumchecks: one for bytecode and one for the matrices, which share their $28$ variables. The aggregation and recursive-statement transcripts use distinct domain labels; each absorbs $n_{\mathrm{rec}}$ before any variable-length sequence, so concatenations cannot be confused. The guest reduces the batch to one bytecode evaluation and a common-point evaluation of each matrix. These three reduced claims, the inner proving-environment digest, and the sub-statements are hashed into the outer VM's two-word public input.

The public \texttt{RecursiveProof::verify} path first verifies that outer VM proof and then evaluates the three fixed polynomials natively at the reduced points. It rejects an empty batch and is the only complete acceptance path. Feeding a \texttt{RecursiveProof} into another aggregation layer (carrying its reduced claims forward rather than checking them natively) is not yet exposed by the current harness.

The matrix sumcheck runs over the $28$-variable polynomial
\[
  P(x) \;=\; A_0(x)\,W_A(x) \;+\; B_0(x)\,W_B(x),
  \qquad
  W_M(x) \;=\; \textstyle\sum_t \gamma_{M,t}\,\eq(r_t,x),
\]
and ends in the reduced claims $\widetilde A_0(r^\ast)$ and $\widetilde B_0(r^\ast)$, at a common point. The in-circuit cost is that of the sumcheck verifier: $28$ cheap rounds. The prover exploits sparsity to avoid the $2^{28}$-sized cube. During the first $14$ rounds (the row variables), each claim enters through the column contraction $m_t[i]=\sum_{j\in\mathrm{row}_i}\eq(r_t^{\mathrm{col}},j)$; for the last $14$ rounds, each matrix is contracted once along its rows, at the row point just fixed. Each contraction is a single addition-only pass over the nonzeros (the entries are Boolean), and all remaining work is dense of size $2^{14}$. Batching $t$ claims per matrix therefore costs $t+1$ passes over its nonzeros plus $O(t\cdot2^{14})$ multiplications, against the $O(2^{28})$ field operations and memory of a dense sumcheck. The bytecode sumcheck is the same, only smaller and with a dense polynomial.
